% =============================================================================
% information.tex
%
% Satu sumber kebenaran untuk data penulis, tesis, institusi, pernyataan AI,
% pengesahan, dan komponen awal lain sesuai PPTA IPB 2026 dan templat
% "format tesis terbaru" (program magister).
%
% Pemuatan: % =============================================================================
% information.tex
%
% Satu sumber kebenaran untuk data penulis, tesis, institusi, pernyataan AI,
% pengesahan, dan komponen awal lain sesuai PPTA IPB 2026 dan templat
% "format tesis terbaru" (program magister).
%
% Pemuatan: % =============================================================================
% information.tex
%
% Satu sumber kebenaran untuk data penulis, tesis, institusi, pernyataan AI,
% pengesahan, dan komponen awal lain sesuai PPTA IPB 2026 dan templat
% "format tesis terbaru" (program magister).
%
% Pemuatan: % =============================================================================
% information.tex
%
% Satu sumber kebenaran untuk data penulis, tesis, institusi, pernyataan AI,
% pengesahan, dan komponen awal lain sesuai PPTA IPB 2026 dan templat
% "format tesis terbaru" (program magister).
%
% Pemuatan: \input{config/information} dari src/main.tex.
% =============================================================================

% --- Opsi Tampilan ---
% Catatan panduan berwarna merah berasal dari templat Word resmi: teks
% instruksi yang harus dihapus pada naskah akhir. Ubah satu baris berikut
% menjadi false agar semua catatan merah hilang sekaligus (judul Daftar
% Tabel/Gambar/Lampiran, Prakata, dan catatan lain yang memakai \ipbcatatan).
% Tata letak tetap sama karena ruang catatan dipertahankan.
\renewcommand{\TampilkanCatatanPanduan}{true}

% --- Data Penulis ---
% Nilai bawaan mengikuti teks placeholder pada templat resmi. Ganti seluruh
% placeholder ini dengan data sebenarnya sebelum naskah diserahkan.
\newcommand{\NamaPenulis}{Nama Penulis}
\newcommand{\NamaPenulisInggris}{Student Name}
\newcommand{\NIM}{NIM}
\newcommand{\TempatLahir}{Kota \ldots{}}
\newcommand{\TanggalLahir}{tanggal bulan tahun}
\newcommand{\EmailPenulis}{alamat.surel@apps.ipb.ac.id}

% --- Data Pembimbing ---
\newcommand{\NamaPembimbingSatu}{Nama Pembimbing 1}
\newcommand{\NIPPembimbingSatu}{}
\newcommand{\NamaPembimbingDua}{Nama Pembimbing 2}
\newcommand{\NIPPembimbingDua}{}
% Pembimbing ketiga bersifat opsional. Set \TampilkanPembimbingTiga menjadi
% true bila komisi pembimbing beranggotakan tiga orang. Nilai ini mengatur
% baris tanda tangan pada halaman pengesahan sekaligus penyebutan nama pada
% Sorotan/Highlights, Ringkasan/Summary, dan Prakata.
\newcommand{\TampilkanPembimbingTiga}{false}
\newcommand{\NamaPembimbingTiga}{Nama Pembimbing 3}

% Daftar pembimbing yang dipakai di dalam kalimat.
\newcommand{\DaftarPembimbing}{%
  \ifdefstring{\TampilkanPembimbingTiga}{true}%
    {\NamaPembimbingSatu, \NamaPembimbingDua, dan \NamaPembimbingTiga}%
    {\NamaPembimbingSatu{} dan \NamaPembimbingDua}%
}
% Versi kapital untuk baris nama pada Sorotan/Highlights dan Ringkasan/Summary.
\newcommand{\DaftarPembimbingKapital}{%
  \ifdefstring{\TampilkanPembimbingTiga}{true}%
    {\MakeUppercase{\NamaPembimbingSatu}, \MakeUppercase{\NamaPembimbingDua},
     dan \MakeUppercase{\NamaPembimbingTiga}}%
    {\MakeUppercase{\NamaPembimbingSatu} dan \MakeUppercase{\NamaPembimbingDua}}%
}
% Versi bahasa Inggris ("and") untuk baris nama pada Highlights/Summary.
\newcommand{\DaftarPembimbingKapitalInggris}{%
  \ifdefstring{\TampilkanPembimbingTiga}{true}%
    {\MakeUppercase{\NamaPembimbingSatu}, \MakeUppercase{\NamaPembimbingDua},
     and \MakeUppercase{\NamaPembimbingTiga}}%
    {\MakeUppercase{\NamaPembimbingSatu} and \MakeUppercase{\NamaPembimbingDua}}%
}

% --- Data Penguji ---
\newcommand{\NamaPengujiSatu}{Nama Lengkap dan Gelar}
\newcommand{\NamaPengujiDua}{Nama Lengkap dan Gelar}

% --- Data Pengesahan ---
\newcommand{\NamaKetuaProgramStudi}{Nama Lengkap dan Gelar}
\newcommand{\NIPKetuaProgramStudi}{}
% Pejabat kedua pada halaman pengesahan: Dekan Fakultas/Sekolah.
\newcommand{\JabatanPejabatKedua}{Dekan Fakultas/Sekolah}
\newcommand{\NamaPejabatKedua}{Nama Lengkap dan Gelar}
\newcommand{\NIPPejabatKedua}{}
% Biarkan kosong agar templat hanya menampilkan label dan ruang pengisian.
\newcommand{\TanggalUjian}{}
\newcommand{\TanggalLulus}{}

% --- Data Tesis ---
\newcommand{\JudulTesis}{%
  Judul Karya Ilmiah Maksimum Tiga Baris, Lima Belas Kata Tidak Termasuk
  Kata Depan dan Kata Sambung%
}
\newcommand{\JudulTesisInggris}{%
  Title of Master's Thesis%
}
\newcommand{\KataKunci}{ditulis dalam bahasa Indonesia, disusun berdasarkan abjad, maksimum lima kata atau frasa}
\newcommand{\Keywords}{written in English, arranged alphabetically, maximum five words or phrases}

% Judul ringkasan visual maksimum enam kata, tidak termasuk kata depan/sambung.
\newcommand{\JudulRingkasanVisual}{%
  Judul Ringkasan Visual Maksimum Enam Kata, Tidak Termasuk Kata Depan dan Kata Sambung}
\newcommand{\JudulRingkasanVisualInggris}{%
  Graphical Summary Title Maximum Six Words, Excluding Prepositions and Conjunctions}
% Isi dengan path berkas gambar untuk mengganti kotak placeholder.
\newcommand{\BerkasRingkasanVisual}{}
\newcommand{\BerkasRingkasanVisualInggris}{}

% --- Deklarasi penggunaan AI ---
% Nilai placeholder mempertahankan pilihan [menggunakan/tidak menggunakan]
% seperti pada DOTX. Ganti dengan true atau false setelah menentukan pilihan.
\newcommand{\MenggunakanAI}{placeholder}
\newcommand{\NamaAlatAI}{[nama alat/layanan]}
\newcommand{\TujuanPenggunaanAI}{[alasan]}

% --- Gelar dan sumber dana ---
\newcommand{\GelarMagister}{magister \ldots{}}
% Kosongkan bila penelitian tidak memperoleh pendanaan khusus.
\newcommand{\SumberDana}{}

% --- Data Institusi ---
\newcommand{\Departemen}{Nama Departemen}
\newcommand{\ProgramStudi}{Nama Program Studi}
\newcommand{\ProgramStudiSingkat}{Nama Program Studi}
\newcommand{\Fakultas}{Fakultas/Sekolah}
\newcommand{\Universitas}{Institut Pertanian Bogor}
\newcommand{\KotaUniversitas}{Bogor}
\newcommand{\Tahun}{20XX}

% --- Data Waktu Pelaksanaan (untuk prakata) ---
\newcommand{\BulanMulaiPenelitian}{\ldots{}}
\newcommand{\TahunMulaiPenelitian}{20XX}
\newcommand{\BulanSelesaiPenelitian}{\ldots{}}
\newcommand{\TahunSelesaiPenelitian}{20XX}
\newcommand{\BulanTahunDokumen}{Bulan Tahun}
 dari src/main.tex.
% =============================================================================

% --- Opsi Tampilan ---
% Catatan panduan berwarna merah berasal dari templat Word resmi: teks
% instruksi yang harus dihapus pada naskah akhir. Ubah satu baris berikut
% menjadi false agar semua catatan merah hilang sekaligus (judul Daftar
% Tabel/Gambar/Lampiran, Prakata, dan catatan lain yang memakai \ipbcatatan).
% Tata letak tetap sama karena ruang catatan dipertahankan.
\renewcommand{\TampilkanCatatanPanduan}{true}

% --- Data Penulis ---
% Nilai bawaan mengikuti teks placeholder pada templat resmi. Ganti seluruh
% placeholder ini dengan data sebenarnya sebelum naskah diserahkan.
\newcommand{\NamaPenulis}{Nama Penulis}
\newcommand{\NamaPenulisInggris}{Student Name}
\newcommand{\NIM}{NIM}
\newcommand{\TempatLahir}{Kota \ldots{}}
\newcommand{\TanggalLahir}{tanggal bulan tahun}
\newcommand{\EmailPenulis}{alamat.surel@apps.ipb.ac.id}

% --- Data Pembimbing ---
\newcommand{\NamaPembimbingSatu}{Nama Pembimbing 1}
\newcommand{\NIPPembimbingSatu}{}
\newcommand{\NamaPembimbingDua}{Nama Pembimbing 2}
\newcommand{\NIPPembimbingDua}{}
% Pembimbing ketiga bersifat opsional. Set \TampilkanPembimbingTiga menjadi
% true bila komisi pembimbing beranggotakan tiga orang. Nilai ini mengatur
% baris tanda tangan pada halaman pengesahan sekaligus penyebutan nama pada
% Sorotan/Highlights, Ringkasan/Summary, dan Prakata.
\newcommand{\TampilkanPembimbingTiga}{false}
\newcommand{\NamaPembimbingTiga}{Nama Pembimbing 3}

% Daftar pembimbing yang dipakai di dalam kalimat.
\newcommand{\DaftarPembimbing}{%
  \ifdefstring{\TampilkanPembimbingTiga}{true}%
    {\NamaPembimbingSatu, \NamaPembimbingDua, dan \NamaPembimbingTiga}%
    {\NamaPembimbingSatu{} dan \NamaPembimbingDua}%
}
% Versi kapital untuk baris nama pada Sorotan/Highlights dan Ringkasan/Summary.
\newcommand{\DaftarPembimbingKapital}{%
  \ifdefstring{\TampilkanPembimbingTiga}{true}%
    {\MakeUppercase{\NamaPembimbingSatu}, \MakeUppercase{\NamaPembimbingDua},
     dan \MakeUppercase{\NamaPembimbingTiga}}%
    {\MakeUppercase{\NamaPembimbingSatu} dan \MakeUppercase{\NamaPembimbingDua}}%
}
% Versi bahasa Inggris ("and") untuk baris nama pada Highlights/Summary.
\newcommand{\DaftarPembimbingKapitalInggris}{%
  \ifdefstring{\TampilkanPembimbingTiga}{true}%
    {\MakeUppercase{\NamaPembimbingSatu}, \MakeUppercase{\NamaPembimbingDua},
     and \MakeUppercase{\NamaPembimbingTiga}}%
    {\MakeUppercase{\NamaPembimbingSatu} and \MakeUppercase{\NamaPembimbingDua}}%
}

% --- Data Penguji ---
\newcommand{\NamaPengujiSatu}{Nama Lengkap dan Gelar}
\newcommand{\NamaPengujiDua}{Nama Lengkap dan Gelar}

% --- Data Pengesahan ---
\newcommand{\NamaKetuaProgramStudi}{Nama Lengkap dan Gelar}
\newcommand{\NIPKetuaProgramStudi}{}
% Pejabat kedua pada halaman pengesahan: Dekan Fakultas/Sekolah.
\newcommand{\JabatanPejabatKedua}{Dekan Fakultas/Sekolah}
\newcommand{\NamaPejabatKedua}{Nama Lengkap dan Gelar}
\newcommand{\NIPPejabatKedua}{}
% Biarkan kosong agar templat hanya menampilkan label dan ruang pengisian.
\newcommand{\TanggalUjian}{}
\newcommand{\TanggalLulus}{}

% --- Data Tesis ---
\newcommand{\JudulTesis}{%
  Judul Karya Ilmiah Maksimum Tiga Baris, Lima Belas Kata Tidak Termasuk
  Kata Depan dan Kata Sambung%
}
\newcommand{\JudulTesisInggris}{%
  Title of Master's Thesis%
}
\newcommand{\KataKunci}{ditulis dalam bahasa Indonesia, disusun berdasarkan abjad, maksimum lima kata atau frasa}
\newcommand{\Keywords}{written in English, arranged alphabetically, maximum five words or phrases}

% Judul ringkasan visual maksimum enam kata, tidak termasuk kata depan/sambung.
\newcommand{\JudulRingkasanVisual}{%
  Judul Ringkasan Visual Maksimum Enam Kata, Tidak Termasuk Kata Depan dan Kata Sambung}
\newcommand{\JudulRingkasanVisualInggris}{%
  Graphical Summary Title Maximum Six Words, Excluding Prepositions and Conjunctions}
% Isi dengan path berkas gambar untuk mengganti kotak placeholder.
\newcommand{\BerkasRingkasanVisual}{}
\newcommand{\BerkasRingkasanVisualInggris}{}

% --- Deklarasi penggunaan AI ---
% Nilai placeholder mempertahankan pilihan [menggunakan/tidak menggunakan]
% seperti pada DOTX. Ganti dengan true atau false setelah menentukan pilihan.
\newcommand{\MenggunakanAI}{placeholder}
\newcommand{\NamaAlatAI}{[nama alat/layanan]}
\newcommand{\TujuanPenggunaanAI}{[alasan]}

% --- Gelar dan sumber dana ---
\newcommand{\GelarMagister}{magister \ldots{}}
% Kosongkan bila penelitian tidak memperoleh pendanaan khusus.
\newcommand{\SumberDana}{}

% --- Data Institusi ---
\newcommand{\Departemen}{Nama Departemen}
\newcommand{\ProgramStudi}{Nama Program Studi}
\newcommand{\ProgramStudiSingkat}{Nama Program Studi}
\newcommand{\Fakultas}{Fakultas/Sekolah}
\newcommand{\Universitas}{Institut Pertanian Bogor}
\newcommand{\KotaUniversitas}{Bogor}
\newcommand{\Tahun}{20XX}

% --- Data Waktu Pelaksanaan (untuk prakata) ---
\newcommand{\BulanMulaiPenelitian}{\ldots{}}
\newcommand{\TahunMulaiPenelitian}{20XX}
\newcommand{\BulanSelesaiPenelitian}{\ldots{}}
\newcommand{\TahunSelesaiPenelitian}{20XX}
\newcommand{\BulanTahunDokumen}{Bulan Tahun}
 dari src/main.tex.
% =============================================================================

% --- Opsi Tampilan ---
% Catatan panduan berwarna merah berasal dari templat Word resmi: teks
% instruksi yang harus dihapus pada naskah akhir. Ubah satu baris berikut
% menjadi false agar semua catatan merah hilang sekaligus (judul Daftar
% Tabel/Gambar/Lampiran, Prakata, dan catatan lain yang memakai \ipbcatatan).
% Tata letak tetap sama karena ruang catatan dipertahankan.
\renewcommand{\TampilkanCatatanPanduan}{true}

% --- Data Penulis ---
% Nilai bawaan mengikuti teks placeholder pada templat resmi. Ganti seluruh
% placeholder ini dengan data sebenarnya sebelum naskah diserahkan.
\newcommand{\NamaPenulis}{Nama Penulis}
\newcommand{\NamaPenulisInggris}{Student Name}
\newcommand{\NIM}{NIM}
\newcommand{\TempatLahir}{Kota \ldots{}}
\newcommand{\TanggalLahir}{tanggal bulan tahun}
\newcommand{\EmailPenulis}{alamat.surel@apps.ipb.ac.id}

% --- Data Pembimbing ---
\newcommand{\NamaPembimbingSatu}{Nama Pembimbing 1}
\newcommand{\NIPPembimbingSatu}{}
\newcommand{\NamaPembimbingDua}{Nama Pembimbing 2}
\newcommand{\NIPPembimbingDua}{}
% Pembimbing ketiga bersifat opsional. Set \TampilkanPembimbingTiga menjadi
% true bila komisi pembimbing beranggotakan tiga orang. Nilai ini mengatur
% baris tanda tangan pada halaman pengesahan sekaligus penyebutan nama pada
% Sorotan/Highlights, Ringkasan/Summary, dan Prakata.
\newcommand{\TampilkanPembimbingTiga}{false}
\newcommand{\NamaPembimbingTiga}{Nama Pembimbing 3}

% Daftar pembimbing yang dipakai di dalam kalimat.
\newcommand{\DaftarPembimbing}{%
  \ifdefstring{\TampilkanPembimbingTiga}{true}%
    {\NamaPembimbingSatu, \NamaPembimbingDua, dan \NamaPembimbingTiga}%
    {\NamaPembimbingSatu{} dan \NamaPembimbingDua}%
}
% Versi kapital untuk baris nama pada Sorotan/Highlights dan Ringkasan/Summary.
\newcommand{\DaftarPembimbingKapital}{%
  \ifdefstring{\TampilkanPembimbingTiga}{true}%
    {\MakeUppercase{\NamaPembimbingSatu}, \MakeUppercase{\NamaPembimbingDua},
     dan \MakeUppercase{\NamaPembimbingTiga}}%
    {\MakeUppercase{\NamaPembimbingSatu} dan \MakeUppercase{\NamaPembimbingDua}}%
}
% Versi bahasa Inggris ("and") untuk baris nama pada Highlights/Summary.
\newcommand{\DaftarPembimbingKapitalInggris}{%
  \ifdefstring{\TampilkanPembimbingTiga}{true}%
    {\MakeUppercase{\NamaPembimbingSatu}, \MakeUppercase{\NamaPembimbingDua},
     and \MakeUppercase{\NamaPembimbingTiga}}%
    {\MakeUppercase{\NamaPembimbingSatu} and \MakeUppercase{\NamaPembimbingDua}}%
}

% --- Data Penguji ---
\newcommand{\NamaPengujiSatu}{Nama Lengkap dan Gelar}
\newcommand{\NamaPengujiDua}{Nama Lengkap dan Gelar}

% --- Data Pengesahan ---
\newcommand{\NamaKetuaProgramStudi}{Nama Lengkap dan Gelar}
\newcommand{\NIPKetuaProgramStudi}{}
% Pejabat kedua pada halaman pengesahan: Dekan Fakultas/Sekolah.
\newcommand{\JabatanPejabatKedua}{Dekan Fakultas/Sekolah}
\newcommand{\NamaPejabatKedua}{Nama Lengkap dan Gelar}
\newcommand{\NIPPejabatKedua}{}
% Biarkan kosong agar templat hanya menampilkan label dan ruang pengisian.
\newcommand{\TanggalUjian}{}
\newcommand{\TanggalLulus}{}

% --- Data Tesis ---
\newcommand{\JudulTesis}{%
  Judul Karya Ilmiah Maksimum Tiga Baris, Lima Belas Kata Tidak Termasuk
  Kata Depan dan Kata Sambung%
}
\newcommand{\JudulTesisInggris}{%
  Title of Master's Thesis%
}
\newcommand{\KataKunci}{ditulis dalam bahasa Indonesia, disusun berdasarkan abjad, maksimum lima kata atau frasa}
\newcommand{\Keywords}{written in English, arranged alphabetically, maximum five words or phrases}

% Judul ringkasan visual maksimum enam kata, tidak termasuk kata depan/sambung.
\newcommand{\JudulRingkasanVisual}{%
  Judul Ringkasan Visual Maksimum Enam Kata, Tidak Termasuk Kata Depan dan Kata Sambung}
\newcommand{\JudulRingkasanVisualInggris}{%
  Graphical Summary Title Maximum Six Words, Excluding Prepositions and Conjunctions}
% Isi dengan path berkas gambar untuk mengganti kotak placeholder.
\newcommand{\BerkasRingkasanVisual}{}
\newcommand{\BerkasRingkasanVisualInggris}{}

% --- Deklarasi penggunaan AI ---
% Nilai placeholder mempertahankan pilihan [menggunakan/tidak menggunakan]
% seperti pada DOTX. Ganti dengan true atau false setelah menentukan pilihan.
\newcommand{\MenggunakanAI}{placeholder}
\newcommand{\NamaAlatAI}{[nama alat/layanan]}
\newcommand{\TujuanPenggunaanAI}{[alasan]}

% --- Gelar dan sumber dana ---
\newcommand{\GelarMagister}{magister \ldots{}}
% Kosongkan bila penelitian tidak memperoleh pendanaan khusus.
\newcommand{\SumberDana}{}

% --- Data Institusi ---
\newcommand{\Departemen}{Nama Departemen}
\newcommand{\ProgramStudi}{Nama Program Studi}
\newcommand{\ProgramStudiSingkat}{Nama Program Studi}
\newcommand{\Fakultas}{Fakultas/Sekolah}
\newcommand{\Universitas}{Institut Pertanian Bogor}
\newcommand{\KotaUniversitas}{Bogor}
\newcommand{\Tahun}{20XX}

% --- Data Waktu Pelaksanaan (untuk prakata) ---
\newcommand{\BulanMulaiPenelitian}{\ldots{}}
\newcommand{\TahunMulaiPenelitian}{20XX}
\newcommand{\BulanSelesaiPenelitian}{\ldots{}}
\newcommand{\TahunSelesaiPenelitian}{20XX}
\newcommand{\BulanTahunDokumen}{Bulan Tahun}
 dari src/main.tex.
% =============================================================================

% --- Opsi Tampilan ---
% Catatan panduan berwarna merah berasal dari templat Word resmi: teks
% instruksi yang harus dihapus pada naskah akhir. Ubah satu baris berikut
% menjadi false agar semua catatan merah hilang sekaligus (judul Daftar
% Tabel/Gambar/Lampiran, Prakata, dan catatan lain yang memakai \ipbcatatan).
% Tata letak tetap sama karena ruang catatan dipertahankan.
\renewcommand{\TampilkanCatatanPanduan}{true}

% --- Data Penulis ---
% Nilai bawaan mengikuti teks placeholder pada templat resmi. Ganti seluruh
% placeholder ini dengan data sebenarnya sebelum naskah diserahkan.
\newcommand{\NamaPenulis}{Nama Penulis}
\newcommand{\NamaPenulisInggris}{Student Name}
\newcommand{\NIM}{NIM}
\newcommand{\TempatLahir}{Kota \ldots{}}
\newcommand{\TanggalLahir}{tanggal bulan tahun}
\newcommand{\EmailPenulis}{alamat.surel@apps.ipb.ac.id}

% --- Data Pembimbing ---
\newcommand{\NamaPembimbingSatu}{Nama Pembimbing 1}
\newcommand{\NIPPembimbingSatu}{}
\newcommand{\NamaPembimbingDua}{Nama Pembimbing 2}
\newcommand{\NIPPembimbingDua}{}
% Pembimbing ketiga bersifat opsional. Set \TampilkanPembimbingTiga menjadi
% true bila komisi pembimbing beranggotakan tiga orang. Nilai ini mengatur
% baris tanda tangan pada halaman pengesahan sekaligus penyebutan nama pada
% Sorotan/Highlights, Ringkasan/Summary, dan Prakata.
\newcommand{\TampilkanPembimbingTiga}{false}
\newcommand{\NamaPembimbingTiga}{Nama Pembimbing 3}

% Daftar pembimbing yang dipakai di dalam kalimat.
\newcommand{\DaftarPembimbing}{%
  \ifdefstring{\TampilkanPembimbingTiga}{true}%
    {\NamaPembimbingSatu, \NamaPembimbingDua, dan \NamaPembimbingTiga}%
    {\NamaPembimbingSatu{} dan \NamaPembimbingDua}%
}
% Versi kapital untuk baris nama pada Sorotan/Highlights dan Ringkasan/Summary.
\newcommand{\DaftarPembimbingKapital}{%
  \ifdefstring{\TampilkanPembimbingTiga}{true}%
    {\MakeUppercase{\NamaPembimbingSatu}, \MakeUppercase{\NamaPembimbingDua},
     dan \MakeUppercase{\NamaPembimbingTiga}}%
    {\MakeUppercase{\NamaPembimbingSatu} dan \MakeUppercase{\NamaPembimbingDua}}%
}
% Versi bahasa Inggris ("and") untuk baris nama pada Highlights/Summary.
\newcommand{\DaftarPembimbingKapitalInggris}{%
  \ifdefstring{\TampilkanPembimbingTiga}{true}%
    {\MakeUppercase{\NamaPembimbingSatu}, \MakeUppercase{\NamaPembimbingDua},
     and \MakeUppercase{\NamaPembimbingTiga}}%
    {\MakeUppercase{\NamaPembimbingSatu} and \MakeUppercase{\NamaPembimbingDua}}%
}

% --- Data Penguji ---
\newcommand{\NamaPengujiSatu}{Nama Lengkap dan Gelar}
\newcommand{\NamaPengujiDua}{Nama Lengkap dan Gelar}

% --- Data Pengesahan ---
\newcommand{\NamaKetuaProgramStudi}{Nama Lengkap dan Gelar}
\newcommand{\NIPKetuaProgramStudi}{}
% Pejabat kedua pada halaman pengesahan: Dekan Fakultas/Sekolah.
\newcommand{\JabatanPejabatKedua}{Dekan Fakultas/Sekolah}
\newcommand{\NamaPejabatKedua}{Nama Lengkap dan Gelar}
\newcommand{\NIPPejabatKedua}{}
% Biarkan kosong agar templat hanya menampilkan label dan ruang pengisian.
\newcommand{\TanggalUjian}{}
\newcommand{\TanggalLulus}{}

% --- Data Tesis ---
\newcommand{\JudulTesis}{%
  Judul Karya Ilmiah Maksimum Tiga Baris, Lima Belas Kata Tidak Termasuk
  Kata Depan dan Kata Sambung%
}
\newcommand{\JudulTesisInggris}{%
  Title of Master's Thesis%
}
\newcommand{\KataKunci}{ditulis dalam bahasa Indonesia, disusun berdasarkan abjad, maksimum lima kata atau frasa}
\newcommand{\Keywords}{written in English, arranged alphabetically, maximum five words or phrases}

% Judul ringkasan visual maksimum enam kata, tidak termasuk kata depan/sambung.
\newcommand{\JudulRingkasanVisual}{%
  Judul Ringkasan Visual Maksimum Enam Kata, Tidak Termasuk Kata Depan dan Kata Sambung}
\newcommand{\JudulRingkasanVisualInggris}{%
  Graphical Summary Title Maximum Six Words, Excluding Prepositions and Conjunctions}
% Isi dengan path berkas gambar untuk mengganti kotak placeholder.
\newcommand{\BerkasRingkasanVisual}{}
\newcommand{\BerkasRingkasanVisualInggris}{}

% --- Deklarasi penggunaan AI ---
% Nilai placeholder mempertahankan pilihan [menggunakan/tidak menggunakan]
% seperti pada DOTX. Ganti dengan true atau false setelah menentukan pilihan.
\newcommand{\MenggunakanAI}{placeholder}
\newcommand{\NamaAlatAI}{[nama alat/layanan]}
\newcommand{\TujuanPenggunaanAI}{[alasan]}

% --- Gelar dan sumber dana ---
\newcommand{\GelarMagister}{magister \ldots{}}
% Kosongkan bila penelitian tidak memperoleh pendanaan khusus.
\newcommand{\SumberDana}{}

% --- Data Institusi ---
\newcommand{\Departemen}{Nama Departemen}
\newcommand{\ProgramStudi}{Nama Program Studi}
\newcommand{\ProgramStudiSingkat}{Nama Program Studi}
\newcommand{\Fakultas}{Fakultas/Sekolah}
\newcommand{\Universitas}{Institut Pertanian Bogor}
\newcommand{\KotaUniversitas}{Bogor}
\newcommand{\Tahun}{20XX}

% --- Data Waktu Pelaksanaan (untuk prakata) ---
\newcommand{\BulanMulaiPenelitian}{\ldots{}}
\newcommand{\TahunMulaiPenelitian}{20XX}
\newcommand{\BulanSelesaiPenelitian}{\ldots{}}
\newcommand{\TahunSelesaiPenelitian}{20XX}
\newcommand{\BulanTahunDokumen}{Bulan Tahun}
